\section{Proof of the regular MaxCut consequence}
\label{sec:regular-maxcut-proof}

We now prove Corollary~\ref{cor:regular-maxcut}. Theorem~\ref{thm:main}
provides one fixed word with SK value close to $\Pstar$, and
Proposition~\ref{prop:tree-transfer} identifies that value with the
large-degree positive-cost tree correlation. The proof first transfers
this correlation to every graph satisfying the stated girth condition.
It then uses local tree structure and the random-regular optimum
asymptotics to obtain the separate successive-limit comparison.

\begin{proof}[Proof of Corollary~\ref{cor:regular-maxcut}]
The normalized negative cost used by Basso et al.~\cite{BFMVZ2022} is
\[
 C_D^-=-\frac1{\sqrt D}\sum_{\{u,v\}\in E}Z_uZ_v
      =\frac2{\sqrt D}H_{\mathrm{MC}}(G)
       -\frac{|E|}{\sqrt D}I.
\]
Thus its cost gate with angle $\gamma_j$ is, up to a scalar phase, the usual
MaxCut cost gate with angle $2\gamma_j/\sqrt D$, with the same sign and
unchanged mixer angle.  Let $r_p^+(D,\gamma,\beta)$ be the distinguished-edge
correlation for the positive tree cost in
Proposition~\ref{prop:tree-transfer}, and put
\[
 \nu_p(D,\gamma,\beta)=\frac{\sqrt D}{2}r_p^+(D,\gamma,\beta).
\]
That proposition gives $\nu_p(D,\gamma,\beta)\to v_p(\gamma,\beta)$ for
each fixed word.

For the observable at an edge, backward cancellation of commuting cost
gates leaves only cost edges having an endpoint at distance at most $p-1$
from the distinguished edge's endpoint set.  An extra edge in this
relevant cone would close a cycle of length at most $2p+1$.  Under the
stated girth condition the cone is therefore the regular tree cone.
Conjugating \emph{only this cone} by $X$ on one bipartition class reverses
the cost sign and the distinguished $Z_uZ_v$, while preserving the mixers
and the initial plus state.  No bipartition of $G$ is used.  Every edge
therefore has negative-cost correlation $-r_p^+(D,\gamma,\beta)$, and
\begin{equation}\label{eq:regular-tree-cut}
 q_{p,G}(2\gamma/\sqrt D,\beta)
 =\frac12+\frac{\nu_p(D,\gamma,\beta)}{\sqrt D}.
\end{equation}
This is the finite-degree edge formula in the full version
of~\cite[Equations~(2.3)--(2.6) and (3.1)]{BFMVZ2022}, in the convention
of Proposition~\ref{prop:tree-transfer}.  Choose the word by
Theorem~\ref{thm:main} with $v_p(\gamma,\beta)\ge\Pstar-\eta/4$.
For all sufficiently large $D$, fixed-word convergence gives
$\nu_p(D,\gamma,\beta)\ge\Pstar-\eta/2$, proving
\eqref{eq:high-girth-maxcut}.

For fixed $D,p$, let $b_{p,D}(G)$ be the fraction of edges whose relevant
cone is not a regular tree cone.  A conservative radius-$(p+1)$
neighborhood test bounds the number of such edges by a constant depending
on $D,p$ times the number of cycles of length at most $2p+3$.
Those fixed-length cycle counts are tight in the uniform simple regular
model by McKay, Wormald, and Wysocka~\cite[Theorem~1]{MWW2004}.
Consequently $b_{p,D}(G_{n,D+1})\to0$ in graph probability and in mean
as $n\to\infty$.  Since each edge contributes a cut probability in $[0,1]$,
the exact value on the other edges gives
\[
 \left|q_{p,G}(2\gamma/\sqrt D,\beta)
 -\frac12-\frac{\nu_p(D,\gamma,\beta)}{\sqrt D}\right|
 \le b_{p,D}(G).
\]
This proves convergence of the conditional quantum expectation to the
tree value at fixed $D,p$; the whole random graph need not have large
girth.

Dembo, Montanari, and Sen~\cite[Theorem~1.6(b), with the convention after
Equation~(1.7)]{DMS2017} identify the optimal leading correction for
the same simple regular model.  After dividing their cut size by $nd/2$
and putting $d=D+1$, their theorem says that
$\sqrt D(\operatorname{OPT}(G_{n,D+1})-1/2)$ converges to $\Pstar$ in
the same successive graph-probability sense; here
$\sqrt{D/(D+1)}\to1$.  Their $\Pstar$ is ours: pairing their ordered
off-diagonal Gaussians and using sign symmetry gives~\eqref{eq:sk-model}
in law, while the diagonal is a mean-zero configuration-independent
shift.  For large $D$ the deterministic quantity $\Pstar-\nu_p(D,\gamma,\beta)$
is at most $\eta/2$.  If $\sqrt D\,\Delta_D>\eta$, the absolute centered-optimum
error or the absolute conditional-expectation error multiplied by
$\sqrt D$ must therefore exceed $\eta/4$.  The first probability vanishes
in the stated successive limit by the cited optimum theorem, and the
second vanishes in the inner $n$ limit by the preceding bound.  This
proves~\eqref{eq:random-regular-maxcut}.
\end{proof}
